\section{Matrix multiplication and exponentiation}

Use matrix exponentiation when a fixed linear transition is applied a huge number of times: linear recurrences, tilings, finite-state DP, and walks of exactly $K$ edges. For Fibonacci, the transition is $\begin{bmatrix}1&1\\1&0\end{bmatrix}$; raising it to power $n-1$ yields $F_n$. For a $k$th-order recurrence, use a $k \times k$ matrix whose first row is the recurrence coefficients and whose remaining rows shift the state.

\begin{lstlisting}
using ll = long long;
const ll MOD = 1e9 + 7;
using Matrix = vector<vector<ll>>;

Matrix multiply(const Matrix& A, const Matrix& B) {
    int n = A.size(), m = B[0].size(), p = B.size();
    Matrix C(n, vector<ll>(m));
    for (int i = 0; i < n; i++) for (int k = 0; k < p; k++) {
        if (A[i][k] == 0) continue; // useful for sparse matrices
        for (int j = 0; j < m; j++)
            C[i][j] = (C[i][j] + A[i][k] * B[k][j]) % MOD;
    }
    return C;
}
Matrix power(Matrix A, long long e) {
    int n = A.size(); Matrix R(n, vector<ll>(n));
    for (int i = 0; i < n; i++) R[i][i] = 1;
    while (e) {
        if (e & 1) R = multiply(R, A);
        A = multiply(A, A); e >>= 1;
    }
    return R;
}
\end{lstlisting}

Matrix multiplication is $O(N^3)$; exponentiation is $O(N^3 \log K)$ with $O(N^2)$ memory. Think of it for huge $n$ (often $10^{18}$), few states, and a transition independent of $n$.

\section{Mo's algorithm}

Mo's algorithm is an offline range-query technique. Reorder queries so the current range \texttt{[L, R]} moves gradually, requiring only \texttt{add(a[i])} and \texttt{remove(a[i])}. It is ideal for static queries such as distinct values, frequency-$K$ counts, sum of squared frequencies, and Powerful Array expressions.

\begin{lstlisting}
struct Query { int l, r, idx; };
int BLOCK;
bool cmp(const Query& a, const Query& b) {
    int x = a.l / BLOCK, y = b.l / BLOCK;
    if (x != y) return x < y;
    return (x & 1) ? a.r > b.r : a.r < b.r;
}

// After sorting queries with cmp:
int L = 0, R = -1;
for (auto q : queries) {
    while (L > q.l) add(a[--L]);
    while (R < q.r) add(a[++R]);
    while (L < q.l) remove(a[L++]);
    while (R > q.r) remove(a[R--]);
    ans[q.idx] = currentAnswer;
}
\end{lstlisting}

Typical complexity is $O((N+Q)\sqrt N \cdot C)$, where $C$ is add/remove cost; memory is $O(N+Q)$. Coordinate-compress large or negative values before using a frequency array. Avoid Mo when queries are online, prefix sums solve the task, or add/remove is expensive. Three-dimensional Mo additionally tracks time and can handle point updates by applying and rolling them back.

\section{Ordered set (PBDS)}

PBDS is a GNU extension for dynamic ordered collections with rank queries. \texttt{find\_by\_order(k)} returns the zero-indexed $k$th smallest; \texttt{order\_of\_key(x)} counts elements strictly smaller than \texttt{x}. Each operation is $O(\log N)$.

\begin{lstlisting}
#include <ext/pb_ds/assoc_container.hpp>
#include <ext/pb_ds/tree_policy.hpp>
using namespace __gnu_pbds;
template<class T> using ordered_set = tree<T, null_type, less<T>,
    rb_tree_tag, tree_order_statistics_node_update>;
template<class T> using ordered_multiset = tree<pair<T,int>, null_type,
    less<pair<T,int>>, rb_tree_tag, tree_order_statistics_node_update>;

ordered_set<int> st;
// *st.find_by_order(k): kth smallest (0-indexed)
// st.order_of_key(x): number of values < x
\end{lstlisting}

PBDS normally rejects duplicates. For a multiset, insert \texttt{\{value, uniqueId\}}. Then \texttt{order\_of\_key(\{X, -1\})} counts values below $X$; a sufficiently large second value counts values at most $X$. Use it for dynamic kth smallest/largest, median, rank, inversion counting, and sliding-window order statistics. It is not standard C++, so do not assume it exists outside GNU contest compilers.

\section{Persistent segment tree}

Persistence keeps old versions after updates. Copy only the root-to-leaf update path and reuse unchanged children; each point update creates $O(\log N)$ nodes and a new root. This supports historical queries and classic kth-smallest-in-subarray queries.

\begin{lstlisting}
struct Node { int l = 0, r = 0; long long sum = 0; };
vector<Node> seg;
int clone(int x) { seg.push_back(seg[x]); return (int)seg.size() - 1; }
int update(int old, int l, int r, int pos, long long val) {
    int cur = clone(old);
    if (l == r) return seg[cur].sum = val, cur;
    int m = (l + r) / 2;
    if (pos <= m) seg[cur].l = update(seg[old].l, l, m, pos, val);
    else seg[cur].r = update(seg[old].r, m + 1, r, pos, val);
    seg[cur].sum = seg[seg[cur].l].sum + seg[seg[cur].r].sum;
    return cur;
}
\end{lstlisting}

Build costs $O(N)$; point updates and range queries cost $O(\log N)$; memory is $O(N + Q\log N)$. For kth smallest in \texttt{[L,R]}, build prefix frequency roots and compare \texttt{root[R]} against \texttt{root[L-1]}. At every node, subtract left-child counts; go left if $K$ fits, otherwise subtract and go right. Coordinate-compress values and map the final rank back to its original value.

\section{Recognition guide}
\begin{itemize}
\item Huge linear recurrence, fixed transition, or exactly $K$ steps: \textbf{matrix exponentiation}.
\item Many static \texttt{[L,R]} queries with easy add/remove: \textbf{Mo's algorithm}.
\item Dynamic ordered values with kth/rank/median queries: \textbf{PBDS}.
\item Old versions, historical queries, or kth smallest in a subarray: \textbf{persistent segment tree}.
\end{itemize}

Useful practice: Fibonacci/Tribonacci and graph walks for matrices; distinct/frequency-squared queries and 3D Mo; dynamic median/rank/inversions for PBDS; and prefix-root kth/count-$\leq X$/versioned updates for persistence. Identify the query structure first, then choose the data structure.
